\section{From Bishop's mixture of experts to our model}

\subsection{Bishop's starting point and the hidden latent variable}

Bishop (\S14.5.3) writes a mixture of experts as
\begin{equation}\label{eq:bishop-moe}
p(y\mid x) = \sum_{k=1}^{K} \pi_k(x)\, p_k(y\mid x).
\end{equation}
(Bishop writes the target as $t$; we write it as $y$.) No latent variable appears in \eqref{eq:bishop-moe}, because it has already been \emph{summed out}. Introduce a one-hot vector $z$ whose component $z_k=1$ means ``expert $k$ generated this observation''. The model then consists of two statements:
\begin{equation}
p(z_k=1\mid x)=\pi_k(x), \qquad p(y\mid x, z_k=1) = p_k(y\mid x).
\end{equation}
Using the exponent trick of Bishop (9.10)--(9.11), the joint distribution is
\begin{equation}
p(y,z\mid x) = p(z\mid x)\,p(y\mid x,z) = \prod_{k=1}^K \big[\pi_k(x)\,p_k(y\mid x)\big]^{z_k}.
\end{equation}
Because $z$ is never observed, we sum it out. Each value of $z$ switches on exactly one factor, so
\begin{equation}
p(y\mid x)=\sum_z p(y,z\mid x)=\sum_{k=1}^K \pi_k(x)\,p_k(y\mid x).
\end{equation}
The latent variable is therefore still present in \eqref{eq:bishop-moe}: $\pi_k$ is its prior and $p_k$ is the distribution of $y$ given its value. It becomes explicit again in two places:
\begin{itemize}
  \item in the \textbf{complete-data likelihood} that EM works with;
  \item in the \textbf{posterior} (the responsibility)
  \begin{equation}
  \gamma_k = p(z_k=1\mid y,x)=\frac{\pi_k(x)\,p_k(y\mid x)}{\sum_j \pi_j(x)\,p_j(y\mid x)}.
  \end{equation}
\end{itemize}

\subsection{Three changes take Bishop's model to ours}

\begin{center}
\renewcommand{\arraystretch}{1.2}
\begin{tabularx}{\textwidth}{@{}lllX@{}}
\toprule
 & Bishop & Ours & Meaning \\
\midrule
Index & $n$ & $(i,t)$ & one observation: stock $i$ on day $t$ \\
Target & $y_n$ & $r_{i,t+1}$ & forward return \\
Expert input & $x_n$ & $x_{i,t}$ & characteristics of the stock \\
Gate input & $x_n$ & $\F_t$ & market information up to $t$ \\
Latent & $z_n$, one per point & $z_{t+1}$, one per date & the regime, shared by all stocks \\
Expert mean & $w_k\T\phi(x)$ & $\mu_k(x_{i,t};\theta_k)$ & linear regression becomes an MLP (Section 3) \\
Noise & $\beta^{-1}$, shared & $s_k^2$, per expert & each regime has its own noise level \\
Gate & $\softmax_k(v\T x_n)$ & HMM probability & Sections 4 and 5 \\
\bottomrule
\end{tabularx}
\end{center}

\textbf{Change 1: the indices.} Bishop has a single index $n$. Our data form a panel, so we use two indices, stock $i$ and day $t$.

\textbf{Change 2: the experts.} Bishop's expert is already a regression whose mean depends on the input. The MLP makes that mean function flexible. Section 3 shows that an MLP is linear regression on learned basis functions.

\textbf{Change 3: the gate.} This change has two parts.
\begin{itemize}
  \item \emph{What the gate looks at.} Our gate sees only market information $\F_t$, so every stock on the same day receives the same weights. The gate answers ``what state is the market in?'', not ``what kind of stock is this?''.
  \item \emph{Where the probability comes from.} Our gate is a probability computed from an explicit model of the regime: a Markov chain over regimes, observed through the market series. It therefore has memory.
\end{itemize}

The result, with every time index written explicitly, is
\begin{equation}\label{eq:our-moe}
\boxed{\;p\big(r_{i,t+1}\mid x_{i,t},\F_t\big)=\sum_{k=1}^K \underbrace{p\big(z_{t+1}=k\mid \F_t\big)}_{\text{HMM gate}}\;\underbrace{\N\big(r_{i,t+1};\,\mu_k(x_{i,t};\theta_k),\,s_k^2\big)}_{\text{expert }k}\;}
\end{equation}
Everything to the right of the conditioning bar is known at the close of day $t$. Only $r_{i,t+1}$ lies in the future. For the study's 5-day target the gate weight is the average of this probability over the window, which reduces to it when $h=1$ (Section~\ref{sec:predfilt}, Q27).

\subsection{The generative story}
For each day:
\begin{enumerate}
  \item the market regime takes one step along its Markov chain;
  \item the market return for that day is drawn from that regime's distribution (this is what the HMM sees);
  \item every stock's return is drawn from the expert of the same regime, with a mean set by the stock's own characteristics.
\end{enumerate}
We never observe step 1. The gate is our best guess of it given the market history, and the mixture averages the experts over that guess.

\textbf{Why an HMM gate rather than Bishop's learned softmax?}
\begin{itemize}
  \item \emph{Persistence.} Regimes last for weeks or months, and the transition matrix builds that memory in.
  \item \emph{Interpretability.} Each state has an economic meaning (for example calm versus stressed), with its own mean and volatility.
  \item \emph{Few parameters.} A Gaussian HMM with $K=2$ has seven parameters (one initial probability, two transition probabilities, two means, two variances), whereas a learned gate is easy to overfit when the signal-to-noise ratio is low.
  \item \emph{Market information enters only through the gate.} Ye \& Borde (2026) found that regime variables helped when they entered through the gate and hurt when they were added to the predictors.
\end{itemize}

\subsection{Information sets as $\sigma$-algebras and filtrations}

``Market information at time $t$'' has a precise meaning in probability theory, and that meaning is what makes ``no look-ahead'' a checkable statement.

\paragraph{The probability space.} All random quantities (regimes $z_t$, market returns $m_t$, characteristics $x_{i,t}$, stock returns $r_{i,t+1}$) live on one probability space $(\Omega,\mathcal A,\mathbb P)$. Here $\Omega$ is the set of possible histories, $\mathcal A$ is a $\sigma$-algebra of events, and $\mathbb P$ is the probability measure.

\paragraph{The gate's information.} The information contained in the market series up to day $t$ is the $\sigma$-algebra it generates,
\begin{equation}
\F_t=\sigma(m_1,\dots,m_t)\subseteq\mathcal A.
\end{equation}
This is the smallest $\sigma$-algebra with respect to which every $m_s$, $s\le t$, is measurable. An event belongs to $\F_t$ exactly when an observer who has seen $m_1,\dots,m_t$ can tell whether it has occurred.

\paragraph{A filtration.} Because $\F_s\subseteq\F_t$ whenever $s\le t$, the family $\mathbb F=(\F_t)_{t\ge1}$ is a \emph{filtration}: information accumulates and is never forgotten. It is the \emph{natural filtration} of the market series. The full information set at the close of $t$ is a larger filtration,
\begin{equation}
\mathcal G_t=\sigma\big(m_s,\ x_{i,s},\ r_{i,s}\ :\ s\le t,\ i\in\mathcal S_{s}\big),\qquad \F_t\subseteq\mathcal G_t.
\end{equation}
The base-case gate is \emph{designed} to use only $\F_t$. The experts' inputs $x_{i,t}$ are $\mathcal G_t$-measurable.

\paragraph{The gate as a conditional probability given a $\sigma$-algebra.}
\begin{equation}
\pi_k(t)=\mathbb P\big(z_{t+1}=k\mid\F_t\big)=\E\big[\mathbf 1\{z_{t+1}=k\}\mid\F_t\big].
\end{equation}
This is a random variable, and it is $\F_t$-measurable. By the Doob--Dynkin lemma it is therefore a measurable function of the observed path: $\pi_k(t)=g_k(m_1,\dots,m_t)$. The forward recursion of Section 5 is the explicit formula for $g_k$.

\paragraph{No look-ahead means adaptedness.} A sequence of forecasts $(\hat r_{i,t+1})_t$ is free of look-ahead exactly when it is \emph{adapted} to the information filtration, i.e.\ $\hat r_{i,t+1}$ is $\mathcal G_t$-measurable for every $t$. The gate must be $\F_t$-measurable. The smoothed probability $\xi_{t\mid T}(k)=\mathbb P(z_t=k\mid\F_T)$ is $\F_T$-measurable but in general \emph{not} $\F_t$-measurable, because it depends on $m_{t+1},\dots,m_T$. That is the formal version of ``smoothed probabilities leak the future'' (Section 8).

\paragraph{Notation.} Wherever these notes write $p(\,\cdot\mid\F_t)$ or $p(\,\cdot\mid m_{1:t})$, they mean the same object: conditioning on $\F_t$ is conditioning on the values of $m_1,\dots,m_t$.

\subsection{Two notational points}

\paragraph{The bar ``$\mid$'' versus the semicolon ``;''.}
\begin{itemize}
  \item $p(a\mid b)$ conditions on a \emph{random variable}: $b$ has a distribution, and $p(a\mid b)=p(a,b)/p(b)$.
  \item $p(a;\theta)$ indexes a family of distributions by a \emph{fixed parameter}. Since $p(\theta)$ is not defined, Bayes' rule cannot be applied to $\theta$.
  \item Frequentist papers such as Jordan \& Jacobs (1994) write $\N(y;\mu,\sigma^2)$. Bishop writes $\mid$ for both cases.
\end{itemize}
A careful version of our expert density puts each quantity where it belongs:
\begin{equation}
p\big(r_{i,t+1}\mid x_{i,t},z_{t+1}=k;\,\theta_k,s_k\big)=\N\big(r_{i,t+1};\,\mu_k(x_{i,t};\theta_k),\,s_k^2\big).
\end{equation}
The rule: anything that will ever be summed over or appear on the left of a probability is a random variable and goes after $\mid$.

\paragraph{The target is a future return, not the next feature vector.}
\begin{itemize}
  \item $x_{i,t+1}$ is tomorrow's \emph{characteristic vector}, whereas the target is a single number, the return $r_{i,t+1}$. In Weigend et al.\ (1995) a series predicts its own next value, which is why they write the target as the next input. In a cross-sectional model the inputs and the target are different variables.
  \item Two dating conventions are in use, and both are fine if applied consistently:
  \begin{itemize}
    \item label by \emph{realisation} date: $y_{i,t}=r_{i,t}$, conditioned on information at $t-1$;
    \item label by \emph{formation} date: $y_{i,t}:=r_{i,t+1}$, conditioned on information at $t$. The code and \path{Model_Derivation_BaseCase.md} use this one.
  \end{itemize}
  \item These notes always write the forward return explicitly: $r_{i,t+1}$ is predicted from $x_{i,t}$, the regime that generates it is $z_{t+1}$, and the market return over the same day is $m_{t+1}$. That way the timing is visible in every equation.
\end{itemize}
